%%%%%%%%%%%%%%%%%%%%%%%%%%%%%%%%%%%%%%%%%%%%%%%%%%%%%%%%%%%%%%%%%%%%%%
% Lớp 10 | Véctơ (chưa xét tọa độ) | Tổng và hiệu của hai véctơ
% Dạng 0H5?2-5: Cực trị hình học
%%%%%%%%%%%%%%%%%%%%%%%%%%%%%%%%%%%%%%%%%%%%%%%%%%%%%%%%%%%%%%%%%%%%%%


% nguồn: De_AI_Toan9_cuoi_hoc_ki_i_1 (nhập 21/11/2026)
\begin{bt}%[0H5V2-5]%[ma:6f4ff99b]
	Cho hình vuông $ABCD$ có cạnh bằng $4$. Gọi $M$ là một điểm di động trên cạnh $BC$. Tìm giá trị nhỏ nhất của biểu thức $P=|\overrightarrow{MA}+\overrightarrow{MD}|$.
	\shortans{8}
	\loigiai{
		Gọi $I$ là trung điểm của cạnh $AD$, ta có $\overrightarrow{MA}+\overrightarrow{MD}=2\overrightarrow{MI}$, suy ra $P=2MI$. \\
Vì $M$ chạy trên cạnh $BC$ song song với $AD$ và khoảng cách giữa hai đường thẳng $BC$ và $AD$ bằng $4$, nên $MI\ge 4$ (dấu bằng xảy ra khi $M$ là trung điểm $BC$). \\
Vậy $P_{\min}=2\cdot 4=8$.
	}
\end{bt}

% nguồn: De_AI_Toan9_cuoi_hoc_ki_i_1 (nhập 21/11/2026)
\begin{bt}%[0H5V2-5]%[ma:d0404b3f]
	Cho tam giác $ABC$ vuông cân tại $A$ có $AB=AC=6$. Điểm $M$ di động trên đường thẳng $BC$. Tìm giá trị nhỏ nhất của bình phương độ dài vectơ $\overrightarrow{u}=\overrightarrow{MB}+\overrightarrow{MC}-\overrightarrow{MA}$.
	\shortans{18}
	\loigiai{
		Dựng hình vuông $ABDC$, ta có $\overrightarrow{AC}=\overrightarrow{BD}$. \\
Biến đổi $\overrightarrow{u}=\overrightarrow{MB}+(\overrightarrow{MC}-\overrightarrow{MA})=\overrightarrow{MB}+\overrightarrow{AC}=\overrightarrow{MB}+\overrightarrow{BD}=\overrightarrow{MD}$. \\
Vì $M$ di động trên $BC$ nên $MD$ nhỏ nhất khi $M$ là hình chiếu vuông góc của $D$ lên $BC$, tức là $M$ trùng tâm hình vuông $ABDC$. \\
Khi đó $MD_{\min}=\dfrac{1}{2}AD=3\sqrt{2}$, suy ra $|\overrightarrow{u}|_{\min}^2=(3\sqrt{2})^2=18$.
	}
\end{bt}

% nguồn: De_AI_Toan9_cuoi_hoc_ki_i_1 (nhập 21/11/2026)
\begin{bt}%[0H5V2-5]%[ma:64d127b8]
	Hai lực $\overrightarrow{F}_1$ và $\overrightarrow{F}_2$ cùng tác dụng vào một chất điểm, có độ lớn lần lượt là $12\text{ N}$ và $16\text{ N}$ với góc giữa hai lực thay đổi tùy ý. Hỏi độ lớn lớn nhất của hợp lực $\overrightarrow{F}=\overrightarrow{F}_1+\overrightarrow{F}_2$ hơn độ lớn nhỏ nhất của hợp lực đó bao nhiêu Niutơn?
	\shortans{24}
	\loigiai{
		Ta có $|F_1-F_2|\le|\overrightarrow{F}_1+\overrightarrow{F}_2|\le F_1+F_2$. \\
Độ lớn lớn nhất của hợp lực là $12+16=28\text{ N}$ (khi hai lực cùng hướng). \\
Độ lớn nhỏ nhất của hợp lực là $|16-12|=4\text{ N}$ (khi hai lực ngược hướng). \\
Hiệu cần tìm là $28-4=24\text{ N}$.
	}
\end{bt}

% nguồn: 2026-11-21_205436_hk_ai_lop10 (nhập 21/11/2026)
\begin{bt}%[0H5C2-5]%[ma:c3d4e5f6]
% nguon: AI tự sinh (Ôn thi HK)
Cho tam giác $ABC$ có $AB=4$, $AC=6$, góc $\widehat{BAC}=60^\circ$. Gọi $M$ là điểm di động trên đường thẳng $BC$. Tìm giá trị nhỏ nhất của biểu thức
$$P = \left|\overrightarrow{MA}+\overrightarrow{MB}\right| + \left|\overrightarrow{MA}+\overrightarrow{MC}\right|.$$
Kết quả làm tròn đến hàng phần mười.
\shortans{5{,}3}
	\loigiai{
Gọi $I$, $J$ lần lượt là trung điểm $AB$, $AC$. Khi đó $\overrightarrow{MA}+\overrightarrow{MB}=2\overrightarrow{MI}$, $\overrightarrow{MA}+\overrightarrow{MC}=2\overrightarrow{MJ}$.
Suy ra $P = 2(MI + MJ)$.
Ta có $MI + MJ \ge IJ$ với mọi $M$, dấu bằng khi $M$ nằm trên đoạn $IJ$ và $M \in BC$, tức $M = IJ \cap BC$.
Tính $IJ$: $IJ$ là đường trung bình của tam giác $ABC$ ứng với cạnh $BC$, nên $IJ = \dfrac{BC}{2}$.
Tính $BC$: $BC^2 = AB^2+AC^2-2\cdot AB\cdot AC\cos60^\circ = 16+36-24=28 \Rightarrow BC=2\sqrt{7}$.
Vậy $IJ = \sqrt{7}$, và $P_{\min} = 2\sqrt{7} \approx 5{,}3$.
}

\end{bt}
